%!TEX root = ./../../springer-mono-format/book.tex

\label{def:set_of_sync_words} 
\IN{synchronizing word} \index{reset word|see{synchronizing word}}
\IN{synchronizing word} 
\IN{set of synchronizing words}
\IN{synchronizing automaton}
\IN{synchronizing state}
An automaton $\semiautomat{A}{Q}{\Sigma}{\delta}$ 
is \emph{synchronizing} (or \emph{resettable})
if there exists a word, called a \emph{reset} or \emph{synchronizing word},
driving it into the same state regardless of its current state. 
Formally, a word $u \in \Sigma^*$ is said to be a \emph{synchronizing} (or \emph{reset}) \emph{word} for~$\mathcal A$
if $\applytrans{\delta}{Q}{u} = \{q\}$ for some $q \in Q$.
In this case, 
we say that $u$ \emph{synchronizes} $\mathcal A$, and that~$q$ is a \emph{synchronizing state} of $\mathcal A$.
The \emph{set of synchronizing words} of~$\mathcal A$
is denoted by $\Syn(\mathcal A) = \{ u \in \Sigma^* \mid \exists q \in Q \colon \delta(Q, u) = \{q\} \}$.
Then $\mathcal A$ is synchronizing if and only if $\Syn(\mathcal A)$ is non-empty.

Synchronizing words
correspond to a sequence of actions (or inputs) that
allow an operator (or observer) to regain
control by resetting the automaton into a well-defined state.
This notion arises in different contexts and applications
in coding theory~\cite{DBLP:books/daglib/0025093,DBLP:journals/iandc/Jurgensen08,KariVolkov2021}, robotics~\cite{Eppstein90,DBLP:journals/algorithmica/Goldberg93,DBLP:conf/focs/Natarajan86,DBLP:conf/lata/Volkov08}, 
symbolic dynamics~\cite{trahtman09}, model checking and verification~\cite{Sandberg04}, and group theory~\cite{ArCamSt2017}.

The introduction of synchronizing automata is often attributed to \Cerny and
his seminal paper~\cite{Cer64}. However, the actual history is more complicated
and the concept can be traced back to earlier work of Laemmel~\cite{Laemmel63}
and Liu~\cite{Liu63}.
For a detailed account on the genesis of this concept,
see the survey~\cite[Sect.~1.3]{Volkov22}.